% Options for packages loaded elsewhere
\PassOptionsToPackage{unicode}{hyperref}
\PassOptionsToPackage{hyphens}{url}
\documentclass[
  ignorenonframetext,
]{beamer}
\newif\ifbibliography
\usepackage{pgfpages}
\setbeamertemplate{caption}[numbered]
\setbeamertemplate{caption label separator}{: }
\setbeamercolor{caption name}{fg=normal text.fg}
\beamertemplatenavigationsymbolsempty
% remove section numbering
\setbeamertemplate{part page}{
  \centering
  \begin{beamercolorbox}[sep=16pt,center]{part title}
    \usebeamerfont{part title}\insertpart\par
  \end{beamercolorbox}
}
\setbeamertemplate{section page}{
  \centering
  \begin{beamercolorbox}[sep=12pt,center]{section title}
    \usebeamerfont{section title}\insertsection\par
  \end{beamercolorbox}
}
\setbeamertemplate{subsection page}{
  \centering
  \begin{beamercolorbox}[sep=8pt,center]{subsection title}
    \usebeamerfont{subsection title}\insertsubsection\par
  \end{beamercolorbox}
}
% Prevent slide breaks in the middle of a paragraph
\widowpenalties 1 10000
\raggedbottom
\AtBeginPart{
  \frame{\partpage}
}
\AtBeginSection{
  \ifbibliography
  \else
    \frame{\sectionpage}
  \fi
}
\AtBeginSubsection{
  \frame{\subsectionpage}
}
\usepackage{iftex}
\ifPDFTeX
  \usepackage[T1]{fontenc}
  \usepackage[utf8]{inputenc}
  \usepackage{textcomp} % provide euro and other symbols
\else % if luatex or xetex
  \usepackage{unicode-math} % this also loads fontspec
  \defaultfontfeatures{Scale=MatchLowercase}
  \defaultfontfeatures[\rmfamily]{Ligatures=TeX,Scale=1}
\fi
\usepackage{lmodern}
\ifPDFTeX\else
  % xetex/luatex font selection
\fi
% Use upquote if available, for straight quotes in verbatim environments
\IfFileExists{upquote.sty}{\usepackage{upquote}}{}
\IfFileExists{microtype.sty}{% use microtype if available
  \usepackage[]{microtype}
  \UseMicrotypeSet[protrusion]{basicmath} % disable protrusion for tt fonts
}{}
\makeatletter
\@ifundefined{KOMAClassName}{% if non-KOMA class
  \IfFileExists{parskip.sty}{%
    \usepackage{parskip}
  }{% else
    \setlength{\parindent}{0pt}
    \setlength{\parskip}{6pt plus 2pt minus 1pt}}
}{% if KOMA class
  \KOMAoptions{parskip=half}}
\makeatother
\setlength{\emergencystretch}{3em} % prevent overfull lines
\providecommand{\tightlist}{%
  \setlength{\itemsep}{0pt}\setlength{\parskip}{0pt}}
% Slide layout defaults for warning-clean scientific decks.
\usepackage{etoolbox}
\IfFileExists{xurl.sty}{\usepackage{xurl}}{}
\IfFileExists{seqsplit.sty}{\usepackage{seqsplit}}{\newcommand{\seqsplit}[1]{#1}}
\protected\def\breakseq#1{\seqsplit{#1}}
\protected\def\breaktt#1{\begingroup\ttfamily\seqsplit{#1}\endgroup}
\setlength{\emergencystretch}{6em}
\tolerance=5000
\hbadness=10000
\hfuzz=1pt
\setlength{\tabcolsep}{2pt}
\AtBeginEnvironment{longtable}{\tiny\renewcommand{\arraystretch}{0.86}\setlength{\tabcolsep}{1pt}}
\AtBeginEnvironment{tabular}{\tiny\renewcommand{\arraystretch}{0.86}\setlength{\tabcolsep}{1pt}}

% Natbib and cross-reference fallbacks — slides don't load natbib
% or cleveref, but combined-PDF manuscript prose may emit these
% commands. The fallback renders citations as a bracketed key list
% and cross-references as detokenized labels so slides don't fail on
% undefined control sequences or raw underscores. \providecommand is
% a no-op if packages load later.
\providecommand{\citep}[1]{[#1]}
\providecommand{\citet}[1]{#1}
\providecommand{\citealp}[1]{#1}
\providecommand{\citeauthor}[1]{#1}
\providecommand{\citeyear}[1]{#1}
\providecommand{\cref}[1]{\texttt{\detokenize{#1}}}
\providecommand{\Cref}[1]{\texttt{\detokenize{#1}}}

\newtheorem{proposition}{Proposition}
\newtheorem{hypothesis}{Hypothesis}
\usepackage{bookmark}
\IfFileExists{xurl.sty}{\usepackage{xurl}}{} % add URL line breaks if available
\urlstyle{same}
% fallback for those not using the hyperref driver hyperxmp:
\makeatletter
\@ifundefined{xmpquote}{\newcommand{\xmpquote}[1]{#1}}{}
\makeatother
\hypersetup{
  hidelinks,
  pdfcreator={LaTeX via pandoc}}

\author{\texorpdfstring{}{}}
\date{}

\begin{document}

\section{Appendix C: Accessibility and
Provenance}\label{appendix-c-accessibility-and-provenance}

\begin{frame}[fragile,allowframebreaks]{Figure Accessibility}
\protect\phantomsection\label{figure-accessibility}
All 18 figures are rendered with a colourblind-safe palette (Wong 2011,
8 colours) and high-contrast labels at publication DPI (300). Each
figure carries a descriptive caption so the visual claims are
recoverable from text alone. The palette avoids red-green colour pairs
that are indistinguishable for deuteranopia and protanopia; when more
than 8 categories are needed, continuous colormaps (\texttt{viridis},
\texttt{plasma}) are used instead of extending the discrete palette.
Font sizes are enforced at \(\geq 16\)pt via a centralized style module,
ensuring readability at both screen and print sizes.
\end{frame}

\begin{frame}[fragile,allowframebreaks]{Provenance Chain}
\protect\phantomsection\label{provenance-chain}
Every reported number is injected from a committed artifact rather than
typed by hand; an unresolved placeholder is a hard error, so the
rendered manuscript can contain no orphaned or stale figures. The
configuration hash and artifact inventory bind the prose to the exact
pipeline run that produced it. The provenance chain is:

\begin{enumerate}
\tightlist
\item
  \breaktt{manuscript/config.yaml} defines the search term, engines,
  taxonomy, and hypotheses
\item
  \breaktt{scripts/01\_literature\_search.py} retrieves records →
  \texttt{corpus.jsonl}
\item
  \breaktt{scripts/02\_meta\_analysis\_pipeline.py} analyses the corpus →
  \texttt{*.json} data files
\item
  \breaktt{scripts/04\_generate\_figures.py} renders figures →
  \texttt{*.png} + \breaktt{figure\_registry.json}
\item
  \breaktt{scripts/05\_inject\_variables.py} computes variables from data
  files → manuscript text
\end{enumerate}

Each figure in \breaktt{figure\_registry.json} records its source data
file, generation parameters, and SHA-256 hash, binding the visual output
to the exact pipeline run. Re-running the pipeline with the same
configuration and seed produces identical data outputs.
\end{frame}

\begin{frame}[fragile,allowframebreaks]{FAIR Data Principles}
\protect\phantomsection\label{fair-data-principles}
The pipeline supports FAIR (Findable, Accessible, Interoperable,
Reusable) data principles:

\begin{itemize}
\tightlist
\item
  \textbf{Findable}: Each record carries persistent identifiers (DOI,
  arXiv ID, OpenAlex ID) that make it findable across databases.
\item
  \textbf{Accessible}: The corpus is stored as plain JSONL, readable by
  any JSON parser; figures are standard PNG files.
\item
  \textbf{Interoperable}: The data model uses standard bibliographic
  fields (title, abstract, authors, DOI, year, venue); nanopublications
  are serialized as RDF/TriG.
\item
  \textbf{Reusable}: The entire pipeline is regenerable from
  \breaktt{manuscript/config.yaml}; re-running with the same
  configuration reproduces identical outputs.
\end{itemize}
\end{frame}

\begin{frame}[allowframebreaks]{Honesty}
\protect\phantomsection\label{honesty}
The default corpus is synthetic and labelled as such; the manuscript
does not present fixture-derived numbers as empirical findings about
modafinil. Live findings require a real retrieval run with regenerated
artifacts --- as produced in this instance, which retrieved 2302 real
records from 7 live engines.
\end{frame}

\end{document}
